% ------------------------------------------------------------------------------
% Este fichero es parte de la plantilla LaTeX para la realización de Proyectos
% Final de Grado, protegido bajo los términos de la licencia GFDL.
% ------------------------------------------------------------------------------

A continuación, se describe la motivación del proyecto \textbf{UniHub} y su alcance. También se incluye un glosario de términos y la organización del resto de la presente documentación.

\section{Motivación}
El catálogo de la oferta universitaria de enseñanza superior en España almacena la información de titulaciones oficiales. Sin embargo, la consulta integrada de universidades públicas y privadas, la verificación de planes de estudio vigentes modificados en el Boletín Oficial del Estado (BOE) y la clasificación de Grados, Másteres y Doctorados resultaba compleja y fragmentada.

La motivación principal de este proyecto es construir \textbf{UniHub}, un sistema informático moderno, robusto y automatizado que rastree, homogeneice, persista y exponga de forma interactiva la oferta universitaria oficial de España, facilitando la visualización de asignaturas, módulos y créditos ECTS cuando la fuente publicada proporcione ese nivel de detalle.

\section{Alcance y Objetivos} 
El alcance del proyecto abarca cuatro fases de ingeniería de software independientes pero integradas:
\begin{enumerate}
    \item \textbf{Desarrollo de un Rastreador (Crawler) en Python (Fase 1):} Capaz de obtener los catálogos oficiales, aplicar filtrado estricto de titulaciones vigentes y autorizadas (descartando extinguidas o no vigentes), clasificar Doctorados, inspeccionar todos los candidatos a PDF del BOE, aplicar resiliencia ante fallos de conexión (pausas de 5 min y cortocircuito a los 15 min), registrar métricas y notificar a la API REST tras completar la recolección.
    \item \textbf{Diseño de Base de Datos y API REST en Python (Fase 2):} Creación de un esquema relacional en PostgreSQL con control de acceso por roles (\texttt{ruct\_api\_user} de solo lectura), pipeline de ingesta ETL, ordenación prioritaria (públicas primero, luego privadas) y construcción de una API REST en FastAPI con métricas físicas de contenedores cgroup.
    \item \textbf{Desarrollo del Portal Web Interactivo (Fase 3):} Aplicación SPA en Vite + React con el sistema de diseño visual de la Universidad de Cádiz (UCA), paginación configurable (5, 10, 20, 50, 100 elementos por página), geolocalización por fórmula de Haversine con distancias integradas en tarjetas, sección de universidades destacadas (top 6 más visitadas), apartado "Sobre Nosotros" (TFG Alejandro Ramos Rodríguez, UCA) y panel de administración aislado en la ruta manual \texttt{/admin}.
    \item \textbf{Sistemas de Contenerización y Orquestación Docker (Fase 4):} Despliegue independiente de 4 contenedores (\texttt{ruct\_db}, \texttt{ruct\_crawler}, \texttt{ruct\_api}, \texttt{ruct\_www}) en una red privada virtualizada con persistencia garantizada de datos.
\end{enumerate}

\section{Glosario de términos} 
\begin{description}
    \item[UniHub:] Nombre oficial de la plataforma web y sistema distribuido objeto de este Trabajo Fin de Grado.
    \item[RUCT:] Registro oficial de origen de datos de Universidades, Centros y Títulos.
    \item[BOE:] Boletín Oficial del Estado.
    \item[ECTS:] European Credit Transfer and Accumulation System (Sistema Europeo de Transferencia y Acumulación de Créditos).
    \item[API REST:] Representational State Transfer Application Programming Interface.
    \item[SPA:] Single Page Application (Aplicación de Página Única).
    \item[UCA:] Universidad de Cádiz.
    \item[ETL:] Extract, Transform and Load (Extracción, Transformación y Carga).
    \item[Web Vitals:] Métricas de rendimiento de la experiencia de usuario en la web (FCP, LCP, TTFB).
    \item[Haversine:] Fórmula matemática para calcular distancias entre dos puntos en una esfera a partir de sus coordenadas geográficas.
\end{description}

\section{Organización del documento}
La presente memoria técnica se estructura formalmente en cuatro partes principales, desglosadas en los siguientes capítulos y anexos:

\begin{itemize}
    \item \textbf{Parte I: Prolegómeno}
    \begin{itemize}
        \item \textbf{Capítulo 1: Introducción.} Expone la motivación de partida, el alcance técnico del sistema distribuido en cuatro fases, el glosario de términos clave y la organización global de la memoria.
        \item \textbf{Capítulo 2: Antecedentes.} Analiza el contexto universitario español, la normativa de ordenación académica (RD 822/2021 y RD 99/2011), las fuentes oficiales de datos (RUCT, BOE, SIIU) y el estado del arte de las plataformas existentes.
        \item \textbf{Capítulo 3: Plan de gestión de proyecto.} Detalla la metodología ágil adoptada, el desglose de dedicación temporal mediante diagrama de Gantt (450 horas de trabajo normativo), el presupuesto económico pormenorizado (recursos humanos, equipamiento, licencias e IA generativa, suministros), la matriz de gestión de riesgos y la estrategia de aseguramiento de la calidad y CI/CD.
    \end{itemize}

    \item \textbf{Parte II: Desarrollo}
    \begin{itemize}
        \item \textbf{Capítulo 4: Requisitos del Sistema.} Especifica los objetivos de negocio y de sistema, junto con el catálogo exhaustivo de requisitos funcionales (RF-01 a RF-17), requisitos no funcionales, requisitos de información y reglas de negocio.
        \item \textbf{Capítulo 5: Análisis del Sistema.} Modela formalmente el sistema mediante diagramas de casos de uso y diagramas de interacción y secuencia para la extracción de planes, sincronización ETL y navegación web.
        \item \textbf{Capítulo 6: Diseño del Sistema.} Describe la arquitectura global bajo el modelo C4 y arquitectura por capas, el diseño relacional de la base de datos en PostgreSQL, la especificación de la API REST y el diseño de la interfaz de usuario con la identidad visual corporativa de la UCA.
        \item \textbf{Capítulo 7: Construcción del Sistema.} Detalla la implementación en código del software, profundizando en la arquitectura modular del crawler (Fase 1) en 7 subpaquetes y 4 pipelines, el motor de ingesta y carga ETL en PostgreSQL (Fase 2), la API REST en FastAPI (Fase 3), el portal web reactivo en React 19 + Vite (Fase 4), y los mecanismos de resiliencia y limpieza segura de datos.
        \item \textbf{Capítulo 8: Pruebas del Sistema.} Expone el plan de verificación integral de calidad, documentando la batería de 38 suites de pruebas unitarias automatizadas, pruebas de integración de API, evaluaciones empíricas contra resoluciones del BOE y benchmarks de rendimiento.
        \item \textbf{Capítulo 9: Despliegue del Sistema.} Describe la infraestructura de contenerización con Docker y Docker Compose, la topología de red aislada, los volúmenes persistentes y las políticas de despliegue y monitorización.
    \end{itemize}

    \item \textbf{Parte III: Epílogo}
    \begin{itemize}
        \item \textbf{Capítulo 10: Conclusiones.} Evalúa el grado de consecución de los objetivos planteados, sintetiza las lecciones aprendidas durante el ciclo de vida del proyecto y propone las líneas de trabajo futuro para la evolución de la plataforma.
        \item \textbf{Capítulo 11: Información sobre Licencia.} Detalla el marco legal y de propiedad intelectual aplicable, adoptando licencias de software libre y código abierto (FOSS).
    \end{itemize}

    \item \textbf{Parte IV: Anexos}
    \begin{itemize}
        \item \textbf{Anexo A: Manual del desarrollador.} Guía paso a paso para la instalación del entorno local, ejecución de suites de pruebas y directrices de contribución.
        \item \textbf{Anexo B: Manual de usuario.} Manual de uso interactivo de la plataforma UniHub para estudiantes, investigadores y administradores.
        \item \textbf{Anexo C: Glosario de Términos y Acrónimos Técnicos.} Diccionario extendido de terminología académica, legal e informática.
        \item \textbf{Anexo D: Esquema Físico DDL de Base de Datos.} Sentencias SQL DDL completas de creación de tablas, restricciones de integridad referencial e índices de optimización.
    \end{itemize}
\end{itemize}
